\section{Exploratory Analysis}
To better understand the structure of the datasets and the internal representations of the model, we conduct an exploratory data analysis (EDA) prior to performing the main experiments. The goal of this analysis is to verify that the datasets capture distinct behaviors, identify potential confounding factors, and observe preliminary geometric structure within the model’s activation space.

We begin by examining the datasets used throughout the experiments. Three datasets are used to represent different behavioral contexts: HelpSteer2 for utility oriented prompts, AdvBench for harmful prompts, and Alpaca for harmless instruction prompts. We first inspect the size of each dataset and sample example prompts so we can verify that they correspond to the intended behavioral categories. HelpSteer prompts typically consist of longer, instruction following tasks designed to evaluate the helpfulness and response quality of AI models. On the other hand, AdvBench prompts are short adversarial requests designed to make the model give unsafe or harmful responses. Alpaca prompts are filled with harmless instructions and general purpose tasks.

Next, we analyze the distribution of prompt lengths across datasets. Using the tokenizer associated with the model, we compute the token length of each prompt and visualize the distribution using normalized histograms. This analysis reveals a significant difference in prompt length between datasets. HelpSteer prompts are substantially longer on average than both AdvBench and Alpaca prompts. This observation shows the difference between the datasets, as HelpSteer contains more complicated instructions while the adversarial prompts in AdvBench tend to be more concise. Identifying this discrepancy is important because prompt length has been shown to have influence over model activations and so acts as a potential confounding factor when analyzing the activation space geometry.

Following this dataset level analysis, we extract internal activations from the model to examine the geometry of its representation space. For each prompt, we run the prompt through the model and capture the final layer hidden state corresponding to the last token. To keep the analysis computationally manageable, we sample a subset of prompts from each dataset (500) and collect their activation vectors. These vectors represent the model’s internal representation of each prompt within the high dimensional activation space.

After collecting the activations, we compute the average activation vector for each dataset. These mean representations are used to construct two directions of interest. The first is a safety direction calculated from the difference between the mean activations of harmful prompts and the mean activations of harmless prompts. The second is a utility direction calculated from the difference between the mean activations of utility prompts and harmless prompts. These vectors are here as preliminary approximations of the directions that will later be calculated using the Difference-of-Means and actSVD methods.

To visualize the overall structure of the activation space, we apply Principal Component Analysis (PCA) to the collected activation vectors and project them into two dimensions. This dimensionality reduction allows us to visualize how prompts from the three datasets are distributed within the representation space. The resulting scatter plot provides a view of how prompts associated with different behaviors cluster in distinct regions of the activation space. 

We also quantify the relationships between these representations by calculating the cosine similarities between the mean activation vectors of each dataset as well as between the derived safety and utility directions. These measurements provide an initial indication of whether the safety related and utility related directions are aligned or orthogonal within the representation space.

Finally, we evaluate the quality of the extracted safety direction by projecting harmful and harmless activation vectors onto this direction. By visualizing the distribution of these projections, we can observe whether the safety direction meaningfully separates harmful prompts from harmless ones. A clear separation between the distributions suggests that the difference of means vector captures a meaningful behavioral dimension related to safety. However, we can see in the figure that there is some overlap in the tails.
